\documentclass[a4paper, 11pt]{article}
\usepackage[utf8]{inputenc}
\usepackage[T1]{fontenc}
\usepackage[italian]{babel}
\usepackage{helvet}
\renewcommand{\familydefault}{\sfdefault}
\usepackage[
  top=2cm,
  bottom=2cm,
  left=2cm,
  right=2cm
]{geometry}

\begin{document}
\pagestyle{empty}
\noindent
\textbf{Titolo italiano}: Confronto di tecniche di steganografia tramite modelli generativi\\
\textbf{Titolo inglese}: Comparing of steganography techniques using generative models\\
\textbf{Candidato}: Matteo Diciotti\\
\textbf{Relatore}: Prof. Daniele Baracchi\\
\textbf{Correlatrice}: Dott.ssa Giulia Bertazzini\\
\textbf{Correlatore}: Prof. Daniele Castellana\\
\textbf{Anno accademico}: 2025-2026\\
\vspace{1ex}
\hrule

\section*{Riassunto}
La steganografia occulta l'esistenza di una comunicazione dentro oggetti multimediali innocui. Mentre la crittografia produce traffico cifrato identificabile, la steganografia mira all'indistinguibilità statistica del flusso comunicativo. Un paradigma recente è la \emph{Steganography without Embedding} (SwE), in cui il messaggio segreto guida la generazione di un nuovo artefatto invece di essere nascosto in un contenuto preesistente. L'esempio principale è Meteor, un algoritmo a chiave simmetrica che sfrutta un Large Language Model per pilotare i token testuali via arithmetic coding, producendo stego-text indistinguibili dall'output del modello.

Questa tesi indaga se un approccio analogo a Meteor sia applicabile al dominio visivo. L'obiettivo non è un sistema pronto all'uso, ma esplorare le condizioni in cui un modello generativo visuale può fungere da sampler steganografico, individuando i fattori che lo rendono possibile o lo limitano.

Il percorso sperimentale parte da VQ-GAN, che abbina un codebook discreto di token visivi a un Transformer autoregressivo (analogia con i modelli linguistici). L'implementazione mostra che encoder e decoder non sono funzioni inverse: circa il 19\% dei token non sopravvive al round-trip attraverso i pixel, impedendo il recupero del messaggio. Il fine-tuning dell'encoder, con decoder e codebook congelati, riduce ma non elimina l'errore, a causa del condizionamento laterale tra posizioni adiacenti. Un esperimento su canale lossless conferma l'equivalenza algoritmica (recupero completo), mostrando che l'ostacolo è nel canale, non nell'algoritmo.

Con Open-MAGVIT2, tokenizer basato sulla \emph{Lookup-Free Quantization} (LFQ), emerge un secondo problema. La LFQ è matematicamente invertibile, ma il canale PNG quantizza i pixel a 8 bit e altera circa il 16.8\% dei token (1.1\% per bit). Poiché ogni token dipende da tutti i precedenti, un singolo errore si propaga a cascata corrompendo le predizioni successive. Il bit error rate end-to-end raggiunge circa il 65\%, aggravato dal disallineamento del flusso binario dovuto al bitrate variabile dell'arithmetic coding.

Per aggirare la propagazione si adotta una strategia alternativa: i bit del messaggio sono scritti direttamente nei piani di bit dei codici latenti di un'immagine esistente (Quantization Index Modulation, QIM), eliminando la dipendenza dal campionamento sequenziale. Con codici Reed-Solomon il recupero è completo, con capacità di 30--50 byte per immagini $256\times 256$ pixel. Questa soluzione è però un watermark robusto ma statisticamente rilevabile, non steganografia indistinguibile: la modulazione dei codici introduce perturbazioni rilevabili, specie alterando più della metà dei 18 piani di bit.

Il risultato centrale è un trade-off tra indistinguibilità statistica (campionamento autoregressivo) e robustezza al canale (modulazione diretta dei codici): le due proprietà sono mutuamente esclusive per questo tipo di sistemi. La trasposizione di SwE al dominio visivo è concettualmente possibile, come mostra l'esperimento su canale lossless, ma è vincolata dall'assenza di un canale privo di rumore tra token e pixel. Proprietà di cui il dominio testuale gode intrinsecamente e su cui si fonda il successo di Meteor.
\end{document}